%==============================================================================
% Research Methods - Group assignment - pitch

%==============================================================================

\documentclass[aspectratio=169]{beamer}

%==============================================================================
% Preamble
%==============================================================================

%---------- Layout ------------------------------------------------------------

\usetheme{hogent}
\usecolortheme{hgwhite}

%---------- Packages ----------------------------------------------------------

\usepackage[dutch]{babel}

\usepackage{booktabs}
\usepackage{multirow,multicol}
\usepackage{eurosym}

%---------- Metadata ----------------------------------------------------------

\title{Pitch onderzoeksvoorstel.}
\subtitle{Interprofessioneel onderwerp -- Research Methods, 25-26, groep 64}
\author{Simon Borghart \and Emiel De Pelsmaeker \and Simon Boey}
\date{\today}

%==============================================================================
% Presentation content
%==============================================================================

\begin{document}

{
    \setbeamertemplate{background}[imgletter]
        {blurred-background-cellphone-coffee-842554.jpg}{A}

    \begin{frame}
        \maketitle
    \end{frame}
}

\begin{frame}
    \frametitle{Inhoud.}

    \tableofcontents
\end{frame}

%==============================================================================
% Deel 1: Probleem
%==============================================================================

\section{Probleem.}

\subsection{Context.}

{

\begin{frame}

    {\huge \textbf{Meer dan honderd sollicitaties op zeven jaar tijd.
    Telkens afgewezen binnen enkele minuten.}}

    \bigskip

    \textbf{Derek Mobley, eiser in Mobley v. Workday}

\end{frame}
}

\begin{frame}
    \frametitle{De trend: AI bij aanwerving.}

    \begin{itemize}
        \item Bedrijven zetten in toenemende mate AI-gedreven platformen
            in om cv's en sollicitaties automatisch te screenen en te
            rangschikken.
        \item Deze systemen verbeteren de efficiëntie van het
            aanwervingsproces.
        \item Maar ze kunnen tegelijk bestaande vooroordelen op basis van
            leeftijd, ras of geslacht reproduceren of versterken.
        \item Sinds mei 2025 en januari 2026 lopen hierover twee
            collectieve rechtszaken in de Verenigde Staten, tegen Workday
            en tegen Eightfold AI.
    \end{itemize}
\end{frame}

\begin{frame}
    \frametitle{Concreet voorbeeld: twee rechtszaken.}

    \small
    \begin{table}
        \begin{tabular}{lp{4.2cm}p{4.2cm}}
        \toprule
                        & \textbf{Mobley v.\ Workday} & \textbf{Kistler v.\ Eightfold AI} \\
        \midrule
        Wet             & ADEA (leeftijdsdiscriminatie) & FCRA (kredietrapportagewet) \\
        Kern            & Platform als ``agent'' van werkgever & Scoring 0--5 en filtering
                          zonder verplichte kennisgeving \\
        Schaal          & Mogelijk miljoenen 40+ sollicitanten & Data van 1+ miljard werknemers \\
        \bottomrule
        \end{tabular}
    \end{table}

    \medskip
    \begin{minipage}{.75\textwidth}
    Eightfold AI wordt onder meer gebruikt door Microsoft, PayPal, Morgan
    Stanley, Starbucks, Chevron en Bayer.
    \end{minipage}
\end{frame}

\subsection{Onderzoek bevestigt het probleem.}

\begin{frame}
    \frametitle{Wat toont onderzoek aan?}

    \small
    \begin{itemize}
        \item Fabris et al.\ (2025, ACM TIST) brengen in een
            multidisciplinair overzicht de technische, juridische en
            organisatorische oorzaken van bias bij algoritmische
            aanwerving in kaart.
        \item Chen (2023) bevestigt via literatuuronderzoek en bevraging
            dat AI-gedreven recrutering leidt tot discriminatoire
            praktijken op basis van onder meer geslacht en ras.
        \item Wen et al.\ (2025) meten met een eigen benchmark
            (\emph{FAIRE}) aanzienlijke raciale en gender-bias in hoe
            grote taalmodellen cv's beoordelen.
        \item Xu et al.\ (2025) tonen aan dat taalmodellen tot 60\% vaker
            cv's shortlisten die door datzelfde model geschreven zijn.
    \end{itemize}
\end{frame}

\begin{frame}
    \frametitle{Probleemstelling.}

    Bedrijven zetten in toenemende mate AI-gedreven platformen in om
    cv's en sollicitaties automatisch te screenen en te rangschikken.
    Deze systemen verbeteren de efficiëntie van het aanwervingsproces,
    maar kunnen tegelijk bestaande vooroordelen op basis van leeftijd,
    ras of geslacht reproduceren of versterken, met reële juridische en
    maatschappelijke gevolgen zoals blijkt uit de lopende rechtszaken
    tegen Workday en Eightfold AI. HR- en compliance-verantwoordelijken
    hebben onvoldoende houvast om bias te herkennen, te meten en te
    beperken, en om te voldoen aan bestaande regelgeving zoals de ADEA
    en de FCRA.
\end{frame}

\subsection{Onderzoeksvraag.}

\begin{frame}
    \frametitle{Onderzoeksvraag.}

    \begin{quote}
        \textbf{Welke combinatie van technische en organisatorische
        maatregelen (bias-audits, transparante scoringscriteria en
        menselijk toezicht op automatische afwijzingen) is het meest
        geschikt voor een werkgever die een AI-aanwervingsplatform zoals
        dat van Workday of Eightfold AI inzet, om discriminatie van
        sollicitanten te beperken zonder de efficiëntiewinst van
        geautomatiseerde screening significant te verminderen?}
    \end{quote}
\end{frame}

\begin{frame}
    \frametitle{Doelgroep.}

    \begin{itemize}
        \item HR- en compliance-verantwoordelijken bij middelgrote tot
            grote organisaties.
        \item Specifiek: teams die een AI-aanwervingsplatform zoals dat
            van Workday of Eightfold AI inzetten of overwegen in te
            zetten.
        \item Verantwoordelijk voor het beperken van juridische en
            reputatierisico's rond geautomatiseerde aanwerving.
        \item Dus niet ``HR-professionals'' in het algemeen, maar
            specifiek de teams die de leverancierskeuze en het
            screeningbeleid bepalen.
    \end{itemize}
\end{frame}

\begin{frame}
    \frametitle{Waarom dit onderwerp?}

    \small
    \begin{itemize}
        \item \textbf{Specific}: drie met naam benoemde maatregelen
            (bias-audits, transparante scoringscriteria, menselijk
            toezicht) binnen \'e\'en concreet probleem.
        \item \textbf{Measurable}: bias-reductie, doorlooptijd en aantal
            (on)terechte afwijzingen zijn kwantificeerbaar, uit
            literatuur \'en uit een eigen PoC-audit.
        \item \textbf{Achievable}: publiek beschikbare datasets,
            open-source auditingtools, geen interne bedrijfsdata nodig.
        \item \textbf{Realistic/Relevant}: recent aangetoond in
            peer-reviewed onderzoek, verankerd bij publiek
            gedocumenteerde bedrijven (Workday, Eightfold AI).
        \item \textbf{Time-bound}: volledige aanpak uitgewerkt binnen
            twaalf weken.
    \end{itemize}
\end{frame}

%==============================================================================
% Deel 2: Aanpak
%==============================================================================

\section{Aanpak.}

\subsection{Deelvragen.}

\begin{frame}
    \frametitle{Deelvragen: probleemdomein.}

    \begin{itemize}
        \item Welke technische oorzaken liggen aan de basis van bias in
            AI-gedreven cv-screening en rangschikkingssystemen zoals die
            van Workday en Eightfold AI?
        \item Welke juridische risico's lopen werkgevers die dergelijke
            systemen inzetten, in het licht van regelgeving zoals de
            ADEA en de FCRA?
        \item In welke mate zijn sollicitanten zich bewust van de rol
            van AI in het aanwervingsproces, en welke impact heeft dit
            op hun vertrouwen in de werkgever?
    \end{itemize}
\end{frame}

\begin{frame}
    \frametitle{Deelvragen: oplossingsdomein.}

    \begin{itemize}
        \item Hoe werkt een bias-audit van een AI-aanwervingssysteem in
            de praktijk, en welke meetmethoden worden hiervoor gebruikt?
        \item Welke vormen van menselijk toezicht op automatische
            afwijzingsbeslissingen worden in de literatuur voorgesteld?
        \item Hoe kan een werkgever transparantie over scoringscriteria
            organiseren zonder de onderliggende modellen volledig
            openbaar te maken?
        \item Hoe kan de impact van deze maatregelen op zowel
            bias-reductie als efficiëntie gemeten en vergeleken worden?
    \end{itemize}
\end{frame}

\subsection{Methodologie.}

\begin{frame}
    \frametitle{Vier fasen, twaalf weken.}

    \small
    \begin{table}
        \begin{tabular}{@{}lcp{5.8cm}@{}}
        \toprule
        \textbf{Fase} & \textbf{Weken} & \textbf{Milestone} \\
        \midrule
        1. Probleemdomein  & 1--3   & Probleembeschrijving + juridisch risicokader \\
        2. Literatuur       & 4--6   & Literatuurgebaseerd afwegingskader \\
        3. Technisch luik   & 7--10  & PoC bias-auditscript op cv-dataset \\
        4. Synthese         & 11--12 & Eindrapport \\
        \bottomrule
        \end{tabular}
    \end{table}
\end{frame}

\begin{frame}
    \begin{itemize}
        \item \textbf{Wat}: een eenvoudig bias-auditscript implementeren
            en toepassen op een publiek beschikbare cv- of
            aanwervingsdataset.
        \item \textbf{Meting}: het verschil in scoring of
            shortlisting-kans tussen kandidaten met vergelijkbare
            kwalificaties maar verschillende gesuggereerde
            demografische kenmerken.
        \item \textbf{Aanpak}: naar analogie met de meetmethode van
            Wen et al.\ (2025).
        \item \textbf{Waarom dit telt}: vult het cesuurcriterium
            ``technische component'' in, naast het theoretische werk.
    \end{itemize}
\end{frame}

\subsection{Resultaat.}

\begin{frame}
    \frametitle{Verwacht resultaat.}

    \begin{itemize}
        \item Een afwegingskader waarmee een werkgever zoals een
            gebruiker van Workday of Eightfold AI bias kan beperken
            zonder de efficiëntie van geautomatiseerde screening te
            verliezen.
        \item Gemeten PoC-resultaten die aantonen hoe detecteerbaar bias
            is met een eenvoudig, haalbaar auditscript.
        \item Geen interne bedrijfsdata nodig: publieke datasets en
            gepubliceerde rechtszaakdocumenten volstaan.
    \end{itemize}
\end{frame}

%==============================================================================
% Tot slot
%==============================================================================

\section{Tot slot.}

\begin{frame}
    \frametitle{Bronnen.}

    \small
    \begin{itemize}
        \item \textbf{Fabris et al.\ (2025)}, ACM TIST -- peer-reviewed
            tijdschriftartikel, multidisciplinair overzicht.
        \item \textbf{Chen (2023)}, Humanities and Social Sciences
            Communications -- peer-reviewed tijdschriftartikel
            (Springer/Nature).
        \item \textbf{Wen et al.\ (2025)}, \textbf{Xu et al.\ (2025)} --
            arXiv-preprints, meting van bias bij cv-screening.
        \item \textbf{Jones Walker LLP (2026)}, \textbf{University of
            Miami Law Review (2026)} -- onafhankelijke juridische
            analyses van beide rechtszaken.
        \item \textbf{EU AI Act (2024)}, \textbf{EEOC (2024)} --
            regelgevend kader, EU en VS.
    \end{itemize}

    \medskip
    \begin{minipage}{.7\textwidth}
    Bewust een mix van brontypes: peer-reviewed publicaties,
    arXiv-preprints \'en onafhankelijke juridische analyses.
    \end{minipage}
\end{frame}

\begin{frame}
    {\Huge \textbf{Vragen?}}

    \bigskip
    Bedankt voor jullie aandacht.
\end{frame}

\end{document}